\chapter{Survival and Exploration}
% \chapter{Survival in the cold}
% \chapter{Survival in the hot}
\begin{multicols}{2}

\section{Basic Concepts}

The survival markers are \textbf{Hunger, Thirst and Exhaustion.} A character must eat to regain energy, drink water to stay hydrated, and rest to alleviate exhaustion.

\textbf{Exploration Turn (ET):} It is convenient to divide a day into blocks of 4 hours because it splits a day into 6 parts. The day has 24 hours: 8 hours of travel, 8 hours of sleep, 4 hours to set up camp and search for food, and 4 hours to fix anything that goes wrong or do something else. Each block of 4 hours is considered an exploration turn. Any action that spends an exploration turn is an exploration action.

\textbf{Point crawl:} The system was designed so that exploration actions could be performed on a point crawl map. Such a map consists of nodes. Each of them can have several events attached. The events can be accessible to characters by being present in the node and achieving good results in exploration tests.

\subsection{Simplified Survival}

Using the complete survival system all the time is unnecessary. Not every situation is a survival situation. Use these simplified rules for long and uneventful travels, and for expeditions that don't need to manage food, water and shelter.

\textbf{Simplified Hunger and Thirst:} consider that characters are always well fed and hydrated. consider that there is enough food and water for the jouney. Charge a reasonable amount for it, if it were bought.  

\textbf{Simplified Resting:} Allow characters to sleep once per day. Heal half of the character's exhaustion and heal $3+health/2$ per sleep.   

\textbf{Simplified Wear:} Traveling causes no wear.

\textbf{Simplified Point Crawl:} keep the node structure and exploration tests for travel. Do not use events meant to reward resources and shelter.

\section{Character Values}

\subsection{Hunger} represents the lack of nutrients in the body. Allowing hunger values to exceed certain levels will apply negative conditions:

\begin{proplist}
    \item = 30 = Becomes Malnourished
    \item 15 - 29 = Becomes Weakened
    \item 0 - 14 = Everything is normal
\end{proplist}

Hunger cannot be increased further than 30. Every hunger point that would be above 30 increases IL by +1 instead, until it eventually leads to death.

Hunger is lowered by eating. Average characters can eat 4 portions per day. Each portion eaten recovers a number of hunger points defined in the portion itself. The nutritional value of a
portion ranges from 1 to 5.

The nutritional value of a portion of food is counted in hunger points. The portions are measured for a character of size 3. Other sizes of characters require different amounts of food to constitute a portion at a rate defined by the VM. Bigger characters need more food for a portion and smaller ones require less.

For roleplaying purposes, a portion of food is roughly 500g.

\subsection{Thirst}

Thirst starts at 0. The higher the character's thirst, the worse it is. Exceeding a certain limit causes the adverse effects below:

\begin{proplist}
    \item 15 = Dead
    \item 12 - 14 = Unconscious + Dehydrated
    \item 8 - 11 = Dehydrated
    \item 4 - 7 = Thirsty
    \item 0 - 3 = Everything is normal
\end{proplist}

Thirst is lowered by drinking water. A character can recover 5 thirst in a single ET if there is enough water available. Every portion of water recovers 1 thirst, unless there is something wrong with it, like being saltwater or infected.

For roleplaying purposes, a portion of water is roughly 0.5L.

\subsection{Rest, Exhaustion and Healing}

\textbf{Exhaustion} is the measure of mental fatigue. It is increased by staying awake and performing mentally demanding tasks, like casting spells or concentrating for long periods. 

Exhaustion starts at zero and increases with every ET, except when resting. The higher it is, the worst the penalties. When the character reaches 12, it has to make a will test every exploration turn to avoid falling unconscious and being forced to short rest. The Will test has a DL equal to character's exhaustion - 5.

\begin{proplist}
    \item > 11 = Confused
    \item 8 - 11 = Exhausted
    \item 4 - 7 = Tired
    \item 0 - 3 = Everything is normal
\end{proplist}

Exhaustion is lowered by long rests and short rests.

\textbf{Sleep } It is an action that costs 2 exploration turns. It heals IL, recovers exhaustion, and can only be done once per day. The character is unconscious during sleep. 

After sleeping, the exhaustion value is set to half its value - exposure. 

\textbf{Healing:} Happens per sleep. The character heals IL equal to $2+(health-exposure)/2$, where health is the health skill and exposure represents the absence of mental and physical stress. Each IL recovered increases hunger and thirst by 1. Healing is not possible if hunger or hydration values are at their maximum value. 

\textbf{Healing Potions:} take their effect during sleep or a short rest. They heal IL, but do not affect the hunger and thirst values. Any healing that is not used during rest is lost.

\textbf{Healing Wounds:} wounds require medicine or magic to heal. Healing this does not increase hunger or thirst.

\textbf{Short Rest:} The character has the option to spend an ET doing nothing, which recovers 1 exhaustion and +1 IL. Every two short rests in a row cost 1 hunger and 1 thirst. Any time a character does not participate in an exploration action, they are doing nothing and short rest automatically.

\textbf{Exposure:} Each terrain has a value that represents how rough it is to spend time there. This represents how physically and mentally stressful the situation is. Generally, the value ranges from 5 to -2 for places that resting is possible. Anything higher than 5 prevents resting, while -2 means you are as comfortable as sleeping in your home's bed. The character can lower exposure by using items.

Optional: \textbf{Infection:} Make a health test every time the character heals through rest. The DL is equal to the 2x Injury penalty + infection exposure. The character takes 4/2/1/0 IL damage depending on the test result.

\textbf{Long Camping:} Sometimes characters need to rest for several days to heal up. In that case, consider that any shelter that was built and any campfires that were made last for several days stay as they are. The exposure value stays constant. Every day, characters heal $6+(health-exposure)/2$. They can continue doing this until water or food runs out. Healing abilities can be used for 1 ET per day per healer. Exhaustion should be zero by the second day for everyone.

\subsection{Wear}

Performing any exploration action that is not resting causes Wear.

\textbf{What cause wear:} The effect is applied to everyone who participated in the exploration action, independent of them rolling for it or not. If they were in the group that executed the exploration action, they receive the respective wear.

\textbf{Standard Wear:} The standard wear is 1 IL + 1 exhaustion per ET, and it is applied at the end of the ET. 

\textbf{Stamina:} Traveling is more exhausting for characters with lower stamina. Use STA to determine extra IL penalties. As a default, STA < 10 = +1 IL per ET. STA < 8 = +2 IL per ET, STA < 6 = +3 IL and so on.

Note: carrying gear and being injured lower STA.

It is possible for the exploration action to have increased wear depending on the terrain. This is explained in the terrain  subsection.

\section{Exploration}

\textbf{Exploration Action:} There are 4 exploration actions: \textbf{Rest}, \textbf{Travel}, \textbf{Search}, and \textbf{Prepare}. The skill test is performed by only one character from the group. Roll individually for each separate group. 

\subsection{Exploring and traveling}

\skill{Exploration}{ Knowledge + Awareness }{ This skill is used to perform Travel and Search actions. Its value is equal to the specific knowledge about the location + the awareness proficiency. 

\textbf{Acquiring Specific knowledge:} Acquiring maps, reaching landmarks, spending time somewhere and asking for information improve local knowledge. These are events unlocked through traveling or in preparation for travel. 

\textbf{Travel:} Tries to move to a different node in the correct direction along the easiest path, avoiding accidents and dangers. The test is made against the DL of the path being attempted. A crit or hit means that no one got lost and the movement is as fast as possible.

\textbf{Search:} Attempts to find events of the desired types and avoid unwanted types without getting lost. The player states what kind of event they are looking for and makes the test. Characters find the easiest event of the type that was searched for and for which their test value was high enough. 

\textbf{Grazes and misses:} If characters dont find any events with search nor the path they were looking for with travel, nor their quarry when tracking, they find a random event in the node. The explanation is in the events subsection.

\textbf{Splitting the party:}The party can be split for any exploration actions. At the end of the ET, the party meets again unless someone moved to another node or did not come back for some reason special reason.
}

\subsection{Making camp}

\textbf{Prepare:} Setting up camp can be time-consuming without proper equipment, but it is essential to avoid exposure and to be able to cook food. Prepare actions include: Campfire, Cooking and Shelter. The skill used is 2x survival knowledge.

\textbf{Campfire:} The campfire allows cooking, provides a source of light and reduces exposure to cold and insects by 1. It requires firewood and some method to light the fire. \textbf{Searching for firewood:} requires characters to perform an exploration action to search for dry firewood. On a crit, they do not spend the exploration action and succeed. On a hit, they succeed. On a graze, they can either give up or spend double the time. On a failure, the action fails. The firewood is enough for 3 ETs.\textbf{Lighting the fire:} is done without a test if some items or spells are present, otherwise, try to use natural materials. Crits perform it in minutes, hits spend an ET, grazes gives the option to spend 2 ETs to succeed or lose one and misses just waste an ET.

\textbf{Cooking:} This test is a survival knowledge test made against the DL of the food and can require equipment to be done, like a campfire and a pot. On a critical success, all food is utilized. On a success, only 75\%. On a graze, 50\%, and on a miss, 25\%.

\textbf{Shelter:} Setting up camp can be time-consuming without proper equipment. Characters can bring their own tent and sleeping bag, which provide a bonus to prevent exposure. If these items are not available, an exploration action to build a natural shelter must be performed. Building a natural shelter takes one ET and improves the exposure for resting by 1 on a hit and by 2 on a crit.

Optionally, roll one random encounter per day when long camping. 

\subsection{Night and Darkness}

Traveling at night significantly reduces the team's visibility range, making exploration more difficult. If the group has a light source or the ability to see in the dark, the DL for search actions is increased by just +5. If the group is without vision, all exploration actions increase by +5 and search increases by +10.

The day can have any number of periods of day and night. This depends on the season and the location on Earth. The standard is 4 periods of day and 2 of night.

\subsection{ How the loop works }

\textbf{Preparing for expeditions:} Having tools to make a campfire, carrying tents and bedrolls, bringing a pot to cook, bringing food and water, and hiring a guide are all measures that should be thought of before starting an expedition. Exploration is a demanding task and can by itself cause as much harm as a battle.

\textbf{Pushing the travel pace:} full day in which characters sleep once and explore 4 times is meant to net a negative IL sum. This is equivalent to pushing the pace.

\textbf{Regular travel pace:} A day where characters explore 3 times, rest once and sleep once is meant to net around 0 IL. This is a good traveling pace.

\textbf{Healing serious injuries:} To heal injuries effectively after tough battles, a long camping is expected to happen. In case of a game with magic, intense use of potions and magic are necessary for quick healing. 

\textbf{Pick your Battles:} Battles are costly things. Even in victory, the injuries received will make players question if it was worth it.

\textbf{Split the party:} Splitting the party allows multiple characters to attempt at a test in the same ET, at the cost of danger of ambushes and accidents for the loners. It is also useful to leave depleted characters resting while others explore. 

\textbf{Baseline Metabolism:} Any action that is not an exploration action nor sleep is considered a short rest. That represents the baseline metabolic rate and it ensures that being still spends resources.

\section{ Point crawl}

\subsection{ Terrain}

Every node in a point crawl map is made out of a terrain, which contains a list of characteristics. The same terrain can be reused for many nodes.

\textbf{Exposure:} Every terrain has a value of exposure. It defines how hard it is to rest in it without finding shelter.

\textbf{Infection Exposure:} Increases the DL of infection tests. 

\textbf{Natural Shelter DL:} Standard difficulty to build a natural shelter. May require tools, like an axe, or shovel.

\textbf{Tracking DL:} Represents how hard it is to track someone. Depends on how long footprints survive, how much the smell sticks, how much sound propagates, the presence of branches to be broken or other things that indicate someone's passage. Default is 0.

\textbf{Visibility:} Defines how far away objects can be spotted. Its value ranges from immediate to horizon. Depends on trees, rocks, fog, etc. Default is shooting range.

\textbf{Wear:} Terrain can have value for extra wear. For example, spending an ET in the desert without shelter always inflicts an extra thirst.

\textbf{Terrain tags:} Determines which actions are effective at reducing exposure in this terrain.
\begin{itemize}
    \item wood: allows looking for firewood. Use natural shelter DL when searching for firewood.
    \item rain: there is a chance of rain, which increases exposure in the absence of shelter.
    \item cold: exposure can be reduced with measures against cold.
    \item heat: exposure can be reduced with measures against heat.
    \item insects: exposure can be reduced with measures against insects.
    \item threats: requires someone to stand guard to avoid or reduce a stress induced increase in exposure.
\end{itemize}

\subsection{Events}

\textbf{Event DL:} The DL of events varies between -5 and 15. This is the difficulty for an event to be found on purpose. A crit over the event DL allows interacting with it and, if the event does not take the whole ET, go for another event for which there was a hit or crit. 

There are 8 types of events:

\begin{proplist}
    \item \textbf{Place:} A structure, village, monument.
    \item \textbf{Encounter:} This is found by the player characters. It can be anything that generates social interaction or combat.
    \item \textbf{Obstacle:} It can be a puzzle, a door, a physical challenge, a natural disaster, etc.
    \item \textbf{Treasure:} It can be any object of value.
    \item \textbf{Wear:} It is an event that increases wear. Can have a chance to happen regardless of the exploration test result.
    \item \textbf{Food:} an event that provides the opportunity to obtain food.
    \item \textbf{Shelter:} finds a place with better exposure than the rest of the node.
    \item \textbf{Water:} an event that provides the opportunity to obtain water.
\end{proplist}

Events can be of multiple types at once. The typing exists so that players can choose to try to search for them. If an event has multiple types, the character must be looking only for one of the types for the event to qualify for the search.

\textbf{Random event:} If none of the desired events or paths were found, put all available events in a numbered list and roll a d6 or d10. Leave numbers blank or add some events multiple times. If the exploration test result is lower than 0, also add paths to other nodes.

\textbf{Random encounter:} Works just like the random events, but happens when characters are making camp. Only wear events and encounter events can be activated.

\textbf{Traveling:} from one node to the next can cost multiple ETs. There is no need for placeholder nodes in the middle of two important ones just to conform to the duration of 1 ET. 

\textbf{Being tracked:} this is rolled separately using the tracker's exploration skill. Trigger a detection check when the trackers approach to determine if they are detected and at what distance.

\textbf{Connection events:} happen when traveling from one node to another. They are triggered automatically by traveling, but often can be avoided if the exploration test was high enough.

% \textbf{Natural Disasters} include: Snowstorm, Sandstorm, Storm, Hail, Earthquake, Eruption, Lightning, Flood, Wildfire, Landslide, Fog, Tornado, strong wind, being covered by insects during sleep, and entering a monster's den.

% \textbf{Animal Encounters:} Here are some ideas for animal encounters: a large fresh carcass is found, but it belongs to a predator. A nest full of eggs is discovered. A large animal is found sleeping in a grove. Something is chasing the group. Hundreds of birds are crossing the sky. A goat is on the edge of a cliff. A whole herd of bison is near the river. The river is full of alligators.

% \textbf{NPC Encounters:} The player characters are not the only ones in the wild. Here are some ideas of who they might encounter: A child is found alone in the forest. A cabin is found in the middle of the path during a snowstorm, but the owner does not want to let anyone in. A merchant is stuck on the road with a broken cart wheel. A werewolf is found camping in the forest and offers shelter in their camp.

Examples of events:

Shelter: DL 2, Reduces the terrain exposure to 0. DL for building natural shelter =0. \\
Encounter: DL 5, A pack of 8 wolves is howling nearby. On a critical success, the party can choose to avoid them. On a success, the party must make a stealth test to avoid the wolves. On a graze, the wolves are aggressive and try to attack. On a miss, the party is caught by surprise.\\
Water: DL 3, A spring of clear water is found. \\
Food: DL 5, a deer is nearby, but it will run away if it gets scared. The deer has enough meat for 80 portions worth 4 and has a DL of 4 for cooking.  
Treasure: A map of the region and 3 GP were found in an abandoned camp. The map increases the effective survival knowledge level in the area by 1 level.
Wear: Traveling causes 1 extra exhaustion to characters with STA lower than 12. Critical successes in exploration reduce the requirement to 10.
Obstacle: A climb is required to continue on the path. Everyone must make a DL 3 climbing test. On a miss, fall down 5m at a DL 3 of balance.

\subsection{Creating the point crawl:}

\textbf{Step 1: Define terrains:} each node has a type of terrain. The first step is to define which terrains are going to be used in the point crawl.

\textbf{Step 2: Define nodes and connections:} Draw the map with each location represented by a node and paths represented by a connection between nodes. Add any points of interest as nodes.

\textbf{Step 3: Travel distances:} For every connection, add a distance measured in ETs. 

\textbf{Define terrain for every node}

\textbf{Add events:} they can be added to nodes or connections. 


\end{multicols}

% \textbf{ \large Water Events}

% Event table:

% \noindent
%

% \begin{tabular}{| c c c |}
% \hline
% Environment & Name & Description \\
% \hline && \\
% Forest  & nothing & nothing \\ 
% Forest  & water puddle & 1L of contaminated water \\ 
% Forest  & stream & unlimited contaminated water \\ 
% Forest & water spring & unlimited clean water \\ 
% \hline
% \end{tabular}



% \textbf{ \large Food Events}

% Event table:

% \noindent
% \begin{tabular}{| c c c |}
% \hline
% Environment & Name & Description \\
% \hline && \\
% Forest & Nothing & Nothing.  \\ 
% Forest & poisoned fruit & 4 servings of 5 Life in fruit. DL 13 health to avoid losing HP instead of recovering.  \\ 
% Forest & dead deer & Deer carcass with 10 servings of 5 Life. DL 15 preparation to make it edible. \\ 
% Forest & wild boar & A wild boar appears. Eating it yields 50 servings of 5 Life. \\ 
% Forest & dry tree & Can bloom with biomancy magic, yielding 20 servings of 10 Life. \\ 
% Forest & fruit tree & Tree with 20 servings of 5 Life in fruit. \\ 
% \hline
% \end{tabular}

% Obstacle Events

% \noindent
% \begin{tabular}{| c c c |}
% \hline
% Die & Name & Description \\
% \hline && \\
% Forest  & Fall & DL 13 balance to avoid falling 6m \\ 
% Forest  & landslide & Rocks fall, DL 10 reflexes to avoid taking 15 bludgeoning damage \\ 
% Forest  & hill & 2km to go around \\ 
% Forest  & strong wind & Exposure DL increases by 5 \\ 
% Forest  & treasure & A backpack with herbs is found  \\ 
% \hline
% \end{tabular}

% Combat Events

% \noindent
% \begin{tabular}{| c c c |}
% \hline
% Die & Name & Description \\
% \hline && \\
% 1-10  & 1 Troll  & 1 Troll in the forest \\ 
% 11-20  & 8 Wolves & 8 Wolves appear and may attack \\ 
% 21-30  & Ambush & 3 people waiting for a trap to activate, detection 15 or 20 slashing damage\\ 
% 31-40  & gelatinous worms & Hole with 3 gelatinous worms \\ 
% 41-50  &  and   & \\ 
% 76-90  & Spiders & 3 Spiders quietly in their web \\ 
% 91-100  &  & Bumping into a smuggler named Wilson \\ 
% \hline
% \end{tabular}

% Shelter Events

% \noindent
% \begin{tabular}{| c c c |}
% \hline
% Die & Name & Description \\
% \hline && \\
% 1-25  & Flood & +2 exposure. The shelter is destroyed within 4 hours. \\ 
% 26-30  & Ant hill & +3 exposure. \\ 
% 31-50  & Nothing & Does not provide protection from anything \\ 
% 51-55  & Cave with bear & exposure DL 5, has a bear with 50 servings of 5 Life \\ 
% 56-75  & Clearing & -3 exposure, -3 to be found \\ 
% 76-100  & Cave & exposure DL 5, +5 DL to track \\ 
% \hline
% \end{tabular}

% Social Events

% \noindent
% \begin{tabular}{| c c c |}
% \hline
% Die & Name & Description \\
% \hline && \\
% 1-25  & Evil wizard & The wizard controls your minds to steal from you. \\ 
% 26-50  & Thief & Surbit tries to stealthily rob you \\ 
% 51-60  & Merchant & merchant and his bodyguard \\ 
% 61-75  & Pilgrims & 2 Pilgrims carrying 30 GP \\ 
% 76-100  & Den & 1 mutant guernicora, 4 wolves, and 2 humans \\ 
% \hline
% \end{tabular}
